\documentclass[../main.tex]{subfiles}
\begin{document}
\chapter{2020-2021学年微积分（一）（下）期末考试}

\section{单项选择题(每小题 3 分，共 18 分)}
\begin{enumerate}
    \item 微分方程 \( y'' - 3y' + 2y = 2x - 2e^x \) 的特解 \( y^* \) 的形式是（\qquad）.

          \begin{tasks}[label=\textbf{\Alph*.}, label-width=1.5em, column-sep=1em](2)
              \task \( y^* = (Ax + B)e^x \)
              \task \( y^* = x(Ax + B)e^x \)
              \task \( y^* = Ax + B + Ce^x \)
              \task \( y^* = Ax + B + Cxe^x \)
          \end{tasks}

    \item 设曲线 \( \begin{cases} x^2 + y^2 + z^2 = 2, \\ x + y + z = 0 \end{cases} \)，点 \( M(1, -1, 0) \)，则在点 \( M \) 处下列说法\textbf{不正确}的是（\qquad）.

          \begin{tasks}[label=\textbf{\Alph*.}, label-width=1.5em, column-sep=1em](2)
              \task 切向量为 \((-2, -2, 4)\)
              \task 切向量为 \((-2, 2, 4)\)
              \task 切线方程为 \( x - 1 = y + 1 = -\dfrac{z}{2} \)
              \task 法平面方程为 \( x + y - 2z = 0 \)
          \end{tasks}

    \item 函数 \( f(x, y) = \sqrt{x^2 + y^2} \) 在点 \( (0, 0) \) 处（\qquad）.

          \begin{tasks}[label=\textbf{\Alph*.}, label-width=1.5em, column-sep=1em](2)
              \task 可微
              \task 偏导数存在
              \task 连续
              \task 不连续
          \end{tasks}

    \item 已知函数 \( f \) 连续，则二次积分 \( \displaystyle \int_0^{\frac{\pi}{2}} \mathrm{d}\theta \displaystyle \int_0^{2\sin\theta} f(r\cos\theta, r\sin\theta) r\,\mathrm{d}r = \)（\qquad）.

          \begin{tasks}[label=\textbf{\Alph*.}, label-width=1.5em, column-sep=1em](1)
              \task \( \displaystyle \int_0^2 \mathrm{d}x \displaystyle \int_0^{1 + \sqrt{1 - x^2}} f(x, y) \mathrm{d}y \)
              \task \( \displaystyle \int_0^2 \mathrm{d}x \displaystyle \int_0^{\sqrt{2x - x^2}} f(x, y) \mathrm{d}y \)
              \task \( \displaystyle \int_0^2 \mathrm{d}y \displaystyle \int_0^{\sqrt{2y - y^2}} f(x, y) \mathrm{d}x \)
              \task \( \displaystyle \int_0^2 \mathrm{d}x \displaystyle \int_0^2 f(x, y) \mathrm{d}y \)
          \end{tasks}

    \item 设 \( \Gamma: \begin{cases} x^2 + y^2 + z^2 = 1, \\ x + y + z = 0 \end{cases} \)，\( I = \displaystyle \oint_{\Gamma} x\,\mathrm{d}s \)，\( J = \displaystyle \oint_{\Gamma} y\,\mathrm{d}s \)，\( K = \displaystyle \oint_{\Gamma} z^2 \mathrm{d}s \).以下说法中正确的是（\qquad）.

          \begin{tasks}[label=\textbf{\Alph*.}, label-width=1.5em, column-sep=1em](2)
              \task \( K = 0 \)
              \task \( I, J, K \) 中有两个等于 0
              \task \( I, J, K \) 都等于 0
              \task \( I, J, K \) 全都不等于 0
          \end{tasks}

    \item 设 \( f(x) \) 是以 \( 2\pi \) 为周期的周期函数且 \( f(x) = \begin{cases} 0, & -\pi \leq x < 0, \\ \pi - x, & 0 \leq x < \pi \end{cases} \)，\( f(x) \) 的傅里叶级数的和函数是 \( S(x) \). 以下说法中正确的是（\qquad）.

          \begin{tasks}[label=\textbf{\Alph*.}, label-width=1.5em, column-sep=1em](2)
              \task \( S(x) \) 处处连续
              \task \( S(x) \equiv f(x) \)
              \task \( S(-1) = 0 \)
              \task \( S(0) = \pi \)
          \end{tasks}
\end{enumerate}

\section{填空题(每小题 4 分，共 16 分)}
\begin{enumerate}
    \setcounter{enumi}{6}
    \item 直线 \( \dfrac{x-2}{2} = \dfrac{y-1}{1} = \dfrac{z-3}{1} \) 与平面 \( x - y + 2z + 4 = 0 \) 的夹角是 \(\underline{\hspace{5em}}\).
          \vspace{1em}

    \item 设 \( P_0(1,1,-1) \)，\( P_1(2,-1,0) \)，则 \( u = x + y^2 + z^3 \) 在 \( P_0 \) 处沿 \( \overrightarrow{P_0P_1} \) 方向的方向导数为 \(\underline{\hspace{5em}}\).
          \vspace{1em}

    \item 若 \( z = z(x,y) \) 由方程 \( \mathrm{e}^{x+2y+3z} + xyz = 1 \) 确定，则 \( z_x(0,0) = \underline{\hspace{5em}}\).
          \vspace{1em}

    \item 级数 \( \sum \limits_{n=0}^{\infty} \dfrac{(-1)^n x^{2n}}{(2n)!} \) (\( -\infty < x < \infty \)) 的和函数为 \(\underline{\hspace{5em}}\).
          \vspace{1em}
\end{enumerate}

\section{基本计算题(每小题 7 分，共 42 分)}

\begin{enumerate}
    \setcounter{enumi}{10}
    \item 求微分方程 \( x \dfrac{\mathrm{d}y}{\mathrm{d}x} = y \ln \dfrac{y}{x} \) 的通解.
          \vspace{10em}

    \item 已知函数 \( z = f(xy, yg(x)) \)，其中 \( f \) 有二阶连续偏导数，\( g \) 可导，求 \( z_x \)，\( z_{xy} \).
          \vspace{10em}

    \item 计算二次积分 \( I = \displaystyle \int_1^3 \mathrm{d}x \displaystyle \int_{x-1}^2 \sin y^2 \mathrm{d}y \).
          \vspace{12em}

    \item 求三重积分 \( I = \displaystyle \iiint \limits_V (x^3 + y^2 + z) \,\mathrm{d}v \)，其中 \( V \) 为 \( x^2 + y^2 + z^2 \leq a^2 \)，\( a > 0 \).
          \vspace{12em}

    \item 求 \( I = \displaystyle \iint \limits_S (z^2 + x) \mathrm{d}y\,\mathrm{d}z - z\,\mathrm{d}x\,\mathrm{d}y \)，其中 \( S \) 是 \( z = \dfrac{1}{2}(x^2 + y^2) \) (\( 0 \leq z \leq 2 \)) 下侧.
          \vspace{12em}

    \item 将 \( f(x) = \arctan x \) 展开为 Maclaurin 级数，并求 \( f^{(20)}(0) \)，\( f^{(21)}(0) \).
          \vspace{12em}
\end{enumerate}

\section{应用题(每小题 7 分，共 14 分)}

\begin{enumerate}
    \setcounter{enumi}{16}
    \item 求 \( f(x, y) = x^4 + y^4 - x^2 - 2xy - y^2 \) 的极值.
          \vspace{10em}

    \item 已知曲线积分 \( \displaystyle \int_L y f(x) \,\mathrm{d}x + \left[f(x) - x^2\right] \,\mathrm{d}y \) 与路径无关，其中 \( f(x) \) 有一阶连续导数，且 \( f(0) = 1 \)，求 \( f(x) \) 和 \( I = \displaystyle \int_{(0,0)}^{(1,1)} y f(x) \,\mathrm{d}x + [f(x) - x^2] \,\mathrm{d}y \) 的值.
          \vspace{10em}
\end{enumerate}

\section{综合题(每小题 5 分，共 10 分)}

\begin{enumerate}
    \setcounter{enumi}{18}
    \item 设 \( D: x^2 + y^2 \leq 1 \)，证明不等式：\( \displaystyle \iint \limits_D \sin \sqrt{(x^2 + y^2)^3} \mathrm{d}x\,\mathrm{d}y \leq \dfrac{2}{5}\pi \).
          \vspace{10em}

    \item 设 \( f(0) = 1 \)，\( f'(0) = 0 \)，\( f''(x) \) 在原点的邻域内有界，证明：\( \alpha > \dfrac{1}{2} \) 时，\( \sum \limits_{n=1}^{\infty} \left( f\left( \dfrac{1}{n^\alpha} \right) - 1 \right) \) 绝对收敛.
          \vspace{10em}
\end{enumerate}
\chapter{2020-2021学年微积分（一）（下）期末考试参考答案}
\section{单项选择题(每小题 3 分，共 18 分)}
\begin{enumerate}
    \item \textbf{Solution}. D.

          方程对应的齐次方程$y''-3y'+2y=0$的特征方程为
          \[
              r^2-3r+2=(r-1)(r-2)=0,
          \]
          有两个不同的实根 \(r_1=1\)，\(r_2=2\).

          对于非齐次项 \(2x\)，因为$\lambda=0$不是特征方程的根，所以可以设特解为 \(y_1^* = Ax + B\).

          对于非齐次项 \(-2e^x\)，因为 \(\lambda=1\) 是特征方程的根，所以可以设特解为 \(y_2^* = Cxe^x\).

          由解的叠加原理，非齐次方程的特解可以设为
          \[
              y^* = y_1^* + y_2^* = Ax + B + Cxe^x.
          \]

    \item \textbf{Solution}. B.

          两曲面的法向量分别为
          \[
              \nabla(x^2+y^2+z^2)=(2x,2y,2z),\quad \nabla(x+y+z)=(1,1,1).
          \]
          在 \(M(1,-1,0)\) 处取值得
          \[
              \boldsymbol{n}_1=(2,-2,0),\quad \boldsymbol{n}_2=(1,1,1).
          \]
          因而曲线的切向量可取作
          \[
              \boldsymbol{\tau} = \boldsymbol{n}_1\times\boldsymbol{n}_2 = (2,-2,0)\times(1,1,1)=(-2,-2,4),
          \]
          所以 A、C、D 都正确，而 B 不是切向方向.

    \item \textbf{Solution}. C.

          因为
          \[
              \lim_{\substack{r\to0,\\\theta\in[0,2\pi]}}r\equiv0,
          \]
          所以函数 \(f(x,y)=\sqrt{x^2+y^2}\) 在 \((0,0)\) 处连续. 又因为
          \[
              \lim_{x\to0}\frac{f(x,0)-f(0,0)}{x}=\lim_{x\to0}\frac{|x|}{x}
          \]
          极限不存在，所以偏导数不存在，更不可微.

    \item \textbf{Solution}. C.

          极坐标区域为
          \[
              0\le \theta\le \frac{\pi}{2},\quad 0\le r\le 2\sin\theta,
          \]
          即第一象限内圆
          \[
              x^2+(y-1)^2\le 1
          \]
          的部分. 改写成直角坐标可得
          \[
              0\le y\le 2,\quad 0\le x\le \sqrt{2y-y^2},
          \]
          故选 C.

    \item \textbf{Solution}. B.

          由对称性可知
          \[
              \begin{gathered}
                  \oint_\Gamma x\,\mathrm{d}s=\oint_\Gamma y\,\mathrm{d}s=\oint_\Gamma z\,\mathrm{d}s=\frac{1}{3}\oint_\Gamma (x+y+z)\,\mathrm{d}s=0,\\
                  \oint_\Gamma z^2\,\mathrm{d}s=\oint_\Gamma y^2\,\mathrm{d}s=\oint_\Gamma x^2\,\mathrm{d}s = \frac{1}{3}\oint_\Gamma (x^2+y^2+z^2)\,\mathrm{d}s=\frac{1}{3}\oint_\Gamma \,\mathrm{d}s\neq0.
              \end{gathered}
          \]
          所以恰有两个积分为零.

    \item \textbf{Solution}. C.

          傅里叶级数在连续点收敛到原函数值，在跳跃点收敛到左右极限平均值. 因为 \(-1\in(-\pi,0)\)，故
          \[
              S(-1)=f(-1)=0.
          \]
          而
          \[
              S(0)=\frac{0+\pi}{2}=\frac{\pi}{2}\neq \pi.
          \]
\end{enumerate}

\section{填空题(每小题 4 分，共 16 分)}
\begin{enumerate}
    \setcounter{enumi}{6}
    \item \textbf{Solution}. $\dfrac{\pi}{6}$.

          直线方向向量为 \((2,1,1)\)，平面法向量为 \((1,-1,2)\). 设夹角为 \(\theta\)，则
          \[
              \sin\theta=\frac{|(2,1,1)\cdot(1,-1,2)|}{\sqrt6\cdot\sqrt6}=\frac12.
          \]
          所以
          \[
              \theta=\frac{\pi}{6}.
          \]

    \item \textbf{Solution}. $0$.

          \[
              \nabla u=(1,2y,3z^2),
          \]
          在 \(P_0(1,1,-1)\) 处为 \((1,2,3)\). 又
          \[
              \overrightarrow{P_0P_1}=(1,-2,1),
          \]
          二者点积为 \(0\)，故方向导数为 \(0\).

    \item \textbf{Solution}. $-\dfrac13$.

          对方程
          \[
              \mathrm{e}^{x+2y+3z}+xyz=1
          \]
          两边关于 \(x\) 求偏导，得
          \[
              \mathrm{e}^{x+2y+3z}(1+3z_x)+yz+xy\,z_x=0.
          \]
          在 \((0,0)\) 处有 \(z=0\)，从而有
          \[
              1+3z_x(0,0)=0,\qquad z_x(0,0)=-\frac13.
          \]

    \item \textbf{Solution}. $\cos x$.

          注意到$\sin x = \sum_{n=0}^{\infty} \frac{(-1)^n x^{2n+1}}{(2n+1)!}$，两端求导得
          \[
              \cos x=\sum_{n=0}^{\infty}\frac{(-1)^n x^{2n}}{(2n)!}.
          \]
\end{enumerate}

\section{基本计算题(每小题 7 分，共 42 分)}
\begin{enumerate}
    \setcounter{enumi}{10}
    \item \textbf{Solution}. 令
          \[
              u=\frac{y}{x},\qquad y=ux.
          \]
          则
          \[
              \frac{\mathrm{d}y}{\mathrm{d}x}=u+x\frac{\mathrm{d}u}{\mathrm{d}x}.
          \]
          原方程化为
          \[
              x\left(u+x\frac{\mathrm{d}u}{\mathrm{d}x}\right)=ux\ln u,
          \]
          即
          \[
              x\frac{\mathrm{d}u}{\mathrm{d}x}=u(\ln u-1).
          \]
          分离变量并积分得
          \[
              \int \frac{\mathrm{d}u}{u(\ln u-1)}=\int \frac{\mathrm{d}x}{x},
          \]
          \[
              \ln\left|\ln u -1\right|=\ln|x| + C,
          \]
          即
          \[
              \ln \frac{y}{x} = Cx + 1.
          \]

    \item \textbf{Solution}. 记
          \[
              u=xy,\qquad v=yg(x),
          \]
          则 \(z=f(u,v)\). 由链式法则
          \[
              z_x=f_1'u_x+f_2'v_x=y(f_1'+g'(x)f_2').
          \]
          再对 \(y\) 求偏导，得
          \[
              z_{xy}=f_1'+g'(x)f_2'+y\left[xf_{11}''+(g(x)+xg'(x))f_{12}''+g'(x)g(x)f_{22}''\right].
          \]

    \item \textbf{Solution}. 原积分区域为
          \[
              1\le x\le 3,\quad x-1\le y\le 2.
          \]
          交换积分次序后可写成
          \[
              0\le y\le 2,\quad 1\le x\le y+1.
          \]
          因而
          \[
              I=\int_0^2\,\mathrm{d}y\int_1^{y+1}\sin y^2\,\mathrm{d}x
              =\int_0^2 y\sin y^2\,\mathrm{d}y.
          \]
          令 \(t=y^2\)，则
          \[
              I=\frac12\int_0^4\sin t\,\mathrm{d}t
              =\frac12(1-\cos4).
          \]

    \item \textbf{Solution}. 由对称性知
          \[
              \iiint\limits_V x^3\,\mathrm{d}v=\iiint\limits_V z\,\mathrm{d}v=0,
          \]
          所以
          \[
              I=\iiint\limits_V y^2\,\mathrm{d}v=\frac{1}{3}\iiint\limits_V (x^2+y^2+z^2)\,\mathrm{d}v.
          \]
          用球坐标代换，
          \[
              I=\frac{1}{3}\int_0^{2\pi}\,\mathrm{d}\theta\int_0^\pi\sin\varphi\,\mathrm{d}\varphi\int_0^a \rho^4\,\mathrm{d}\rho
              =\frac{4\pi a^5}{15}.
          \]

    \item \textbf{Solution}. 将 \(S\) 与上盖面
          \[
              S_1:\ z=2,\quad x^2+y^2\le 4
          \]
          补成闭曲面，其所围区域记为 \(V\). 因 \(S\) 取下侧，为闭曲面的外侧方向，因而 \(S_1\) 取上侧.

          由 Gauss 公式可得
          \[
              \begin{aligned}
                  I = & \oiint\limits_{S+S_1} \left(z^2+x\right)\,\mathrm{d}y\,\mathrm{d}z - z\,\mathrm{d}x\,\mathrm{d}y
                  - \iint\limits_{S_1} \left(z^2+x\right)\,\mathrm{d}y\,\mathrm{d}z - z\,\mathrm{d}x\,\mathrm{d}y            \\
                  =   & \iiint\limits_V \left(1-1\right)\,\mathrm{d}v + 2\iint\limits_{x^2+y^2\le 4}\mathrm{d}x\,\mathrm{d}y \\
                  =   & 8\pi.
              \end{aligned}
          \]

    \item \textbf{Solution}. 由
          \[
              \frac{1}{1+x^2}=\sum_{n=0}^{\infty}(-1)^n x^{2n}\qquad (|x|<1)
          \]
          两端积分，得
          \[
              \arctan x=\sum_{n=0}^{\infty}\frac{(-1)^n}{2n+1}x^{2n+1},\qquad |x|\le 1.
          \]
          比较 Maclaurin 展开系数可知
          \[
              f^{(20)}(0)=0,\qquad f^{(21)}(0)=20!.
          \]
\end{enumerate}

\section{应用题(每小题 7 分，共 14 分)}
\begin{enumerate}
    \setcounter{enumi}{16}
    \item \textbf{Solution}. 令$\begin{cases}
                  f_x=4x^3-2x-2y=0, \\
                  f_y=4y^3-2x-2y=0
              \end{cases}$ 可得驻点
          \[
              (0,0),(1,1),(-1,-1).
          \]
          再由
          \[
              f_{xx}=12x^2-2,\quad f_{yy}=12y^2-2,\quad f_{xy}=-2
          \]
          可知在$(1,1)$和$(-1,-1)$处$f_{xx}f_{yy}-f_{xy}^2>0$，所以$(1,1)$和$(-1,-1)$是$f$的极小值点，极小值为
          \[
              f(1,1)=f(-1,-1)=-2.
          \]

          对于驻点$(0,0)$，令$y=x$，则$f(x,y)=2x^4-4x^2=2x^2(x^2-2)$，当 \(|x|>0\) 充分小时有$f(x,y)<0$；

          再令$y=-x$，则$f(x,y)=2x^4$，当 \(|x|>0\) 充分小时有$f(x,y)>0$，

          这意味着在 $(0,0)$ 的任意邻域内，$f(x,y)$既能取到正值也能取到负值，

          因此 $f(0,0)=0$ 既不是极小值也不是极大值.

    \item \textbf{Solution}. 记
          \[
              P(x,y)=yf(x),\qquad Q(x,y)=f(x)-x^2.
          \]
          由路径无关得
          \[
              P_y=Q_x,
          \]
          即
          \[
              f'(x)-f(x)=2x.
          \]
          解得
          \[
              f(x)=C\mathrm{e}^x-2x-2.
          \]
          再由 \(f(0)=1\) 得 \(C=3\)，故
          \[
              f(x)=3\mathrm{e}^x-2x-2.
          \]
          取折线路径 \((0,0)\to(1,0)\to(1,1)\)，则
          \[
              \begin{aligned}
                  I = & \int_{(0,0)}^{(1,0)} yf(x)\,\mathrm{d}x + \int_{(1,0)}^{(1,1)} \left(3\mathrm{e}^x-2x-2-x^2\right)\,\mathrm{d}y \\
                  =   & (3\mathrm{e}-5)\int_0^1 \,\mathrm{d}y                                                                           \\
                  =   & 3\mathrm{e}-5.
              \end{aligned}
          \]
\end{enumerate}

\section{综合题(每小题 5 分，共 10 分)}
\begin{enumerate}
    \setcounter{enumi}{18}
    \item \textbf{Proof}. 用极坐标代换，
          \[
              \iint\limits_D \sin \sqrt{(x^2+y^2)^3}\,\mathrm{d}x\,\mathrm{d}y
              =\int_0^{2\pi}\,\mathrm{d}\theta\int_0^1 \sin(r^3)\,r\,\mathrm{d}r
              =2\pi\int_0^1 r\sin(r^3)\,\mathrm{d}r.
          \]
          由 \(\sin t\le t\)（\(t\ge 0\)）得
          \[
              2\pi\int_0^1 r\sin(r^3)\,\mathrm{d}r
              \le 2\pi\int_0^1 r\cdot r^3\,\mathrm{d}r
              =2\pi\int_0^1 r^4\,\mathrm{d}r
              =\frac{2\pi}{5}.
          \]
          于是原不等式成立.

    \item \textbf{Proof}. 因 \(f''(x)\) 在 \((-1,1)\) 内有界，设
          \[
              |f''(x)|\le M\qquad (|x|<1).
          \]
          当 \(n\) 充分大时，\(x_n=\dfrac{1}{n^\alpha}\in(-1,1)\). 由 Taylor 公式可得
          \[
              f(x_n)=f(0)+f'(0)x_n+\frac{f''(\theta_n x_n)}{2}x_n^2,
          \]
          其中 \(0<\theta_n<1\). 利用 \(f(0)=1,\ f'(0)=0\) 得
          \[
              f\left(\frac{1}{n^\alpha}\right)-1
              =\frac{f''\left(\frac{\theta_n}{n^\alpha}\right)}{2n^{2\alpha}}.
          \]
          所以
          \[
              \left|f\left(\frac{1}{n^\alpha}\right)-1\right|
              \le \frac{M}{2n^{2\alpha}}.
          \]
          当 \(\alpha>\dfrac12\) 时，级数$\sum_{n=1}^{\infty}\frac{1}{n^{2\alpha}}$收敛，于是由比较判别法知
          \[
              \sum_{n=1}^{\infty}\left|f\left(\frac{1}{n^\alpha}\right)-1\right|
          \]
          收敛，即原级数绝对收敛.
\end{enumerate}
\end{document}






